% Generated draft: requires mathematical and source review.
\documentclass[11pt]{article}
\usepackage[margin=25mm]{geometry}
\usepackage{fontspec}
\setmainfont{Latin Modern Roman}
\usepackage{amsmath,amssymb,mathtools,unicode-math}
\setmathfont{Latin Modern Math}
\usepackage{enumitem,needspace,xurl,hyperref}
\setlength{\parindent}{0pt}
\setlength{\parskip}{5pt}
\setlength{\emergencystretch}{3em}
\newcommand{\N}{\mathbb{N}}
\newcommand{\Z}{\mathbb{Z}}
\newcommand{\Q}{\mathbb{Q}}
\newcommand{\R}{\mathbb{R}}
\newcommand{\C}{\mathbb{C}}
\newcommand{\F}{\mathbb{F}}
\newcommand{\A}{\mathbb{A}}
\newcommand{\PP}{\mathbb{P}}
\newcommand{\E}{\mathbb{E}}
\newcommand{\eps}{\varepsilon}
\newcommand{\dd}{\,\mathrm{d}}
\begin{document}
\section*{UnsolvedMath ID: 30005396}
\textbf{Source problem code:} catalogue entry
\subsection{Definitions and mathematical statement}
\textbf{Scope of the supplied passage.} The record for permanent UnsolvedMath ID \(30005396\) contains a local nonvanishing assertion and a theorem giving cases of an unnamed conjecture. It does not state that conjecture. Consequently, a complete unresolved mathematical problem cannot be recovered from this record. The following preserves the supplied assertions, with the missing hypotheses explicitly identified.

\textbf{Local notation and recorded assertion.} Write \(H_v=H(F_{0,v})\), and let \(\pi_v^\vee\) denote the contragredient of the representation \(\pi_v\) appearing in the excerpt. Here \(\boxtimes\) denotes the external tensor product, and \(\operatorname{Hom}_{H_v}(\pi_v,\C)\) denotes the space of \(H_v\)-invariant linear functionals, with the trivial action on \(\C\); the required representation category is not specified. The excerpt describes an Ichino--Ikeda canonical invariant linear functional
\[
\alpha_v\in\operatorname{Hom}_{H_v\times H_v}
  (\pi_v\boxtimes\pi_v^\vee,\C)
\]
by the matrix-coefficient integral
\[
\alpha_v(\phi_v,\phi_v^\vee)
 =\int_{H(F_{0,v})}
   \langle\pi_v(h)\phi_v,\phi_v^\vee\rangle_v\,\dd h,
\qquad
\phi_v\in\pi_v,\quad\phi_v^\vee\in\pi_v^\vee,
\]
and records the equivalence
\[
\alpha_v\ne0
\quad\Longleftrightarrow\quad
\operatorname{Hom}_{H(F_{0,v})}(\pi_v,\C)\ne0.
\]
The integral is a reconstruction of the damaged source notation. The ambient groups, the embedding giving the action of \(H_v\), the pairing, the Haar measure, convergence or regularization, and any normalization factors are not supplied. These formulas therefore record the excerpt's assertion subject to its missing ambient hypotheses; they do not assert it for arbitrary groups and representations.

\textbf{Recorded proved cases, with the conclusion missing.} The passage labels the following assertion ``Theorem 1'': the unnamed conjecture holds under the following numbered conditions, presented together:
\begin{enumerate}
\item \(p>2n\).
\item The extension \(F/F_0\) is unramified everywhere, and all \(2\)-adic places are split in \(F\).
\item Every place \(v\mid p\) is split in \(F\), and the corresponding \(\pi_v\) are unramified and ordinary.
\item The representation \(\pi\) is stable cuspidal. For the local factorization
\[
\pi_v=\pi_{n,v}\boxtimes\pi_{n+1,v},
\]
the following additional conditions are imposed:
\begin{itemize}
\item If \(v\) is split, at least one of \(\pi_{n,v}\) and \(\pi_{n+1,v}\) is unramified.
\item If \(v\) is inert, each of these two factors is either unramified or almost unramified. The excerpt explains ``almost unramified'' as admitting a nonzero invariant vector under a parahoric subgroup defined by a vertex lattice of type \(1\).
\end{itemize}
\end{enumerate}
The ambient meanings and hypotheses for \(F,F_0,p,n,H,\pi\), the precise notions of stability and ordinarity, and the range of places in condition (4) are absent. The excerpt also does not specify the scope of ``unramified everywhere.'' The supplied curation associates the unnamed conjecture with a \(p\)-adic height conjecture in arithmetic Gan--Gross--Prasad theory, but gives neither its height identity nor its definitions and quantifiers. Thus the displayed hypotheses cannot be turned into a complete theorem or an unresolved conjecture from the supplied evidence alone.
\subsection{Short English statement}
The supplied passage records a local nonvanishing criterion and lists hypotheses under which an unstated \(p\)-adic height conjecture is said to hold. It supplies neither the conjecture itself nor the full ambient setup. Its supplied ``partially solved'' label is dated curation, while its statement-verification field marks the problem formulation as unrecoverable. This statement-only draft requires human mathematical and source review.
\subsection{Sources}
\begin{itemize}
\item Local UnsolvedMath registry: \texttt{data/unsolvedmath/problems.json}, permanent ID \(30005396\), source problem code \texttt{catalogue entry}, title ``Local Nonvanishing in Arithmetic Gan--Gross--Prasad Theory.'' The supplied record belongs to the Oberwolfach Reports Open Problems collection and gives proposed year \(2023\); no proposer is identified.
\item Supplied report citation: \emph{Arithmetic of Shimura Varieties} (2023), Oberwolfach Reports, DOI \texttt{10.4171/owr/2023/5}, pp.\ 261--326. \url{https://doi.org/10.4171/owr/2023/5}. The background describes these reports as documenting mathematical workshops and their open-problem discussions. No precise page within this range is supplied for the excerpt.
\item Supplied literature source: Disegni and Zhang, \emph{GGP cycles and p-adic L-functions}. \url{https://arxiv.org/abs/2410.08401}. This is the attribution supplied by the record; the excerpt's reference ``[1]'' is not explicitly identified by a bibliography.
\item Supplied literature source: \emph{Oberwolfach Report 5/2023}. \url{https://doi.org/10.4171/owr/2023/5}.
\item Literature assessment dated exactly \texttt{2026-08-22}: the supplied status is ``partially solved.'' The supplied assessment says that the recorded theorem gives proved cases of the \(p\)-adic height conjecture, including split ordinary and specified stable cuspidal situations, and does not give a general proof. This is supplied dated curation derived from the OWR record, not an independently verified assessment; its time-sensitive status and cited claims require independent verification.
\item The supplied statement status is ``unrecoverable.'' Its statement-verification field reports that the official report gives a local nonvanishing fact followed by proved cases, without formulating an unresolved problem. This verification description is supplied metadata, not a claim of source consultation in preparing this draft.
\end{itemize}
\end{document}
